% ============================================================
% CAPITOLO 2 — L'Ambiente di Sviluppo
% ============================================================
\chapter{L'Ambiente di Sviluppo: VS~Code, Copilot e Claude~Code}

\section*{Cosa Imparerai}
Alla fine di questo capitolo avrai:
\begin{itemize}
  \item VS Code installato e configurato per lo sviluppo 0-code
  \item GitHub Copilot attivo e funzionante (oppure Claude Code)
  \item Generato il tuo primo frammento di codice con l'IA
  \item Compreso le tre modalità di interazione con l'IA in VS~Code
\end{itemize}

% ---------------------------------------------------------
\section{Installazione di Visual Studio Code}

Visual Studio Code (VS~Code) è l'editor di codice gratuito di Microsoft. È lo strumento standard dell'industria e, soprattutto, è l'ambiente dove GitHub Copilot e Claude Code operano al massimo delle loro capacità.

\begin{versionbox}
Le istruzioni di questo capitolo sono verificate con VS~Code 1.96+, GitHub Copilot Chat v0.25+ e Claude Code v1.x (Q1~2025). Le interfacce di questi strumenti evolvono rapidamente. Se un pulsante o un menu non corrisponde esattamente a quanto descritto, consulta il repository companion del libro per le istruzioni aggiornate: i concetti e il metodo restano identici.
\end{versionbox}

\begin{praticabox}[title={Pratica --- Installazione VS Code}]
\textbf{Windows:}
\begin{enumerate}
  \item Vai su \url{https://code.visualstudio.com}
  \item Scarica l'installer per Windows
  \item Esegui l'installer. Seleziona tutte le opzioni di integrazione
\end{enumerate}

\textbf{macOS:}
\begin{lstlisting}[language=bash]
brew install --cask visual-studio-code
\end{lstlisting}

\textbf{Linux (Ubuntu/Debian):}
\begin{lstlisting}[language=bash]
sudo snap install code --classic
\end{lstlisting}
\end{praticabox}

\begin{checkpointbox}
Apri VS~Code. Dovresti vedere la schermata di benvenuto. Se sì, procedi.
\end{checkpointbox}

% ---------------------------------------------------------
\section{Le Estensioni Essenziali}

VS~Code da solo è un editor di testo avanzato. La sua potenza viene dalle estensioni.

\begin{praticabox}[title={Pratica --- Installazione estensioni}]
Apri il pannello estensioni (\texttt{Ctrl+Shift+X} su Windows, \texttt{Cmd+Shift+X} su macOS) e installa:
\end{praticabox}

\begin{center}
\begin{tabularx}{\textwidth}{l l X}
\toprule
\textbf{Estensione} & \textbf{ID} & \textbf{A cosa serve} \\
\midrule
GitHub Copilot & \texttt{GitHub.copilot} & Il motore IA principale \\
Copilot Chat & \texttt{GitHub.copilot-chat} & Chat e agent mode \\
ESLint & \texttt{dbaeumer.vscode-eslint} & Linting JavaScript/TypeScript \\
Prettier & \texttt{esbenp.prettier-vscode} & Formattazione automatica \\
Thunder Client & \texttt{rangav.vscode-thunder-client} & Test API REST \\
Flutter & \texttt{Dart-Code.flutter} & Sviluppo Flutter (Part~IV) \\
Prisma & \texttt{Prisma.prisma} & Evidenziazione schema database \\
\bottomrule
\end{tabularx}
\end{center}

\begin{tipbox}
Non installare subito Flutter e Prisma. Li installerai quando arriverai ai capitoli corrispondenti. Per ora ti bastano Copilot e le estensioni JavaScript.
\end{tipbox}

% ---------------------------------------------------------
\section{Configurazione di GitHub Copilot}

GitHub Copilot è lo strumento IA di GitHub (Microsoft) integrato in VS~Code. Utilizza modelli come GPT-4o e Claude Sonnet per generare codice, rispondere a domande e modificare file.

\begin{praticabox}[title={Pratica --- Attivazione Copilot}]
\begin{enumerate}
  \item \textbf{Account GitHub}: Se non ne hai uno, crealo su \url{https://github.com}.
  \item \textbf{Abbonamento Copilot}:
    \begin{itemize}
      \item \textbf{Copilot Free}: Limitato (2000 completamenti + 50 messaggi chat/mese) ma sufficiente per i primi 3--4 capitoli
      \item \textbf{Copilot Individual}: 10\$/mese --- sufficiente per l'intero libro
      \item \textbf{Copilot Pro}: 39\$/mese --- Agent Mode illimitato
    \end{itemize}
  \item \textbf{Login in VS~Code}:
    \begin{itemize}
      \item Apri VS~Code
      \item Cerca ``GitHub: Sign In'' nella Command Palette (\texttt{Ctrl+Shift+P})
      \item Autorizza l'accesso con il tuo account GitHub
    \end{itemize}
  \item \textbf{Verifica}: Apri un file \texttt{.py} o \texttt{.js} e inizia a digitare. Se vedi suggerimenti in grigio, Copilot è attivo.
\end{enumerate}
\end{praticabox}

\begin{costbox}
Con Copilot Individual (10\$/mese), il costo totale per seguire tutti i 16 capitoli è di circa 20--30\,\euro, a seconda del tuo ritmo. Con Claude Code (pay-per-use), il costo è paragonabile: circa 15--40\,\euro{} per i token consumati durante tutto il percorso. In entrambi i casi, è meno di un singolo corso online o di un'ora di consulenza con uno sviluppatore.
\end{costbox}

\subsection{Le tre modalità di Copilot}

Copilot in VS~Code ha tre modalità di interazione, ognuna con un livello di autonomia crescente:

\begin{center}
\begin{tabularx}{\textwidth}{l l X X}
\toprule
\textbf{Modalità} & \textbf{Come si attiva} & \textbf{Cosa fa} & \textbf{Quando usarla} \\
\midrule
Completamento & Automatico & Suggerisce completamenti riga per riga & Scrittura manuale \\
Chat & Pannello laterale & Conversazione libera, genera file & Spiegazioni, snippet \\
Agent Mode & Nella chat $\rightarrow$ ``Agent'' & Autonomia completa & \textbf{Questa è la modalità 0-code} \\
\bottomrule
\end{tabularx}
\end{center}

\begin{versionbox}
L'interfaccia della chat di Copilot (posizione dei pulsanti, scorciatoie, nome della modalità Agent) cambia con gli aggiornamenti di VS~Code. Se non trovi un elemento, cerca nella Command Palette (\texttt{Ctrl+Shift+P}): digita ``Copilot'' per vedere tutte le opzioni disponibili. Quello che conta è la \textit{funzionalità}: la modalità che consente a Copilot di creare file, eseguire comandi e iterare autonomamente.
\end{versionbox}

\begin{warnbox}
L'Agent Mode è la modalità che userai per il 90\% del lavoro in questo libro. Copilot in Agent Mode può creare file, eseguire comandi nel terminale, installare pacchetti e iterare autonomamente sui risultati. È l'equivalente di avere un junior developer che lavora al tuo fianco.
\end{warnbox}

\begin{versionbox}
I modelli disponibili cambiano frequentemente. Al momento della stesura, i modelli raccomandati sono Claude Sonnet~4 (migliore aderenza al contesto), GPT-4o (veloce e versatile) e Gemini~2.5 Pro (ottima finestra di contesto). La regola generale: scegli il modello più recente della famiglia Claude o GPT. Tutti i concetti e le tecniche del libro funzionano con qualsiasi modello sufficientemente avanzato --- non sei vincolato a uno specifico.
\end{versionbox}

% ---------------------------------------------------------
\section{Alternativa: Claude Code da Terminale}

Se preferisci lavorare da terminale, o se hai bisogno di un'autonomia ancora maggiore, puoi usare \textbf{Claude Code}\index{Claude Code} --- lo strumento a riga di comando di Anthropic~\cite{anthropic:claude-code}.

\begin{praticabox}[title={Pratica --- Installazione Claude Code}]
\begin{lstlisting}[language=bash]
# Richiede Node.js 18+
npm install -g @anthropic-ai/claude-code

# Primo avvio e autenticazione
claude

# Navigare nella cartella del progetto e avviare
cd mio-progetto
claude
\end{lstlisting}
\end{praticabox}

\begin{center}
\begin{tabularx}{\textwidth}{l X X}
\toprule
\textbf{Aspetto} & \textbf{Copilot Agent Mode} & \textbf{Claude Code} \\
\midrule
Interfaccia & GUI integrata in VS~Code & Terminale \\
Curva di apprendimento & Bassa & Media \\
Autonomia & Alta & Molto alta \\
Ideale per & Progetti nuovi, lavoro visuale & Progetti complessi, multi-file \\
Costo & 10\$/mese (Individual) & Pay-per-use \\
\bottomrule
\end{tabularx}
\end{center}

\begin{tipbox}
Per questo libro, tutti gli esempi sono presentati con Copilot Agent Mode in VS~Code. Se usi Claude Code, i concetti sono identici --- cambia solo l'interfaccia. Ogni \texttt{\_CONTEXT.md} che scriverai funziona con entrambi.
\end{tipbox}

% ---------------------------------------------------------
\section{Oltre VS Code: Gli IDE AI-Nativi del 2026}

Il panorama degli strumenti di sviluppo si è trasformato. Nel 2026 sono emersi IDE progettati \textbf{dalle fondamenta} attorno all'orchestrazione dell'IA, non come estensioni applicate a un editor esistente.

\begin{center}
\begin{small}
\begin{tabularx}{\textwidth}{l X X X}
\toprule
\textbf{Caratteristica} & \textbf{VS Code + Copilot} & \textbf{Cursor} & \textbf{Windsurf} \\
\midrule
Architettura & Estensione per IDE & Fork AI-nativo di VS Code & IDE standalone RAG \\
Costo base & 10\$/mese & 20\$/mese & 20\$/mese \\
Punto di forza & Ecosistema, estensioni, Git nativo & Composer: modifica multi-file coordinata & RAG profondo su monorepo \\
Contesto & File allegati e workspace & Indicizzazione semantica (200K token) & Estrazione dinamica RAG \\
Ideale per & Progetti didattici, team eterogenei & Architetture complesse, refactoring & Prototipazione rapida, enterprise \\
\bottomrule
\end{tabularx}
\end{small}
\end{center}

\begin{notebox}
Anche Claude Code si è evoluto significativamente: opera con finestre di contesto fino a 1~milione di token e supporta sessioni remote persistenti (\textit{Remote Control}), permettendo di delegare task massivi a un agente che opera in background su macchine virtuali, anche per ore.
\end{notebox}

\subsection{Perché questo libro usa VS Code}

La scelta di VS~Code come ambiente principale non è casuale:

\begin{itemize}
  \item \textbf{Universalità}: gratuito, multipiattaforma, il più adottato al mondo
  \item \textbf{Bassa barriera d'ingresso}: 10\$/mese per Copilot, contro i 20\$ dei concorrenti
  \item \textbf{Ecosistema}: Git nativo, migliaia di estensioni, supporto first-party di Microsoft
  \item \textbf{Trasferibilità}: le competenze di Context Engineering (\texttt{\_CONTEXT.md}, \texttt{SKILL.md}) funzionano \textbf{identicamente} su Cursor, Windsurf o qualsiasi altro IDE
\end{itemize}

\begin{warnbox}
Gli strumenti cambiano velocemente: Cursor potrebbe integrarsi in VS~Code, Windsurf potrebbe cambiare modello di business, potrebbero emergere alternative superiori. Il \textbf{metodo} che impari in questo libro --- Context Engineering, ADLC, Confidence Tagging --- è indipendente dallo strumento e sopravvive a qualsiasi migrazione.
\end{warnbox}

% ---------------------------------------------------------
\section{Il Tuo Primo Codice Generato dall'IA}

È ora di sporcarsi le mani.

\begin{praticabox}[title={Pratica --- Generare una funzione Python}]
\begin{enumerate}
  \item \textbf{Crea una cartella di lavoro:} Apri VS~Code e apri una cartella vuota (es.\ \texttt{test-0code})
  \item \textbf{Apri Copilot Chat in Agent Mode} (dalla Command Palette: cerca ``Copilot Chat'', poi seleziona la modalità ``Agent'')
  \item \textbf{Scrivi questa richiesta:}
\end{enumerate}
\begin{lstlisting}
Crea un file Python chiamato calcolatrice.py che implementi
una calcolatrice da riga di comando. Deve supportare le 4
operazioni base (+, -, *, /). L'utente inserisce due numeri
e l'operazione da terminale. Gestisci la divisione per zero
con un messaggio di errore chiaro.
Aggiungi un file di test con pytest.
\end{lstlisting}
\begin{enumerate}
  \setcounter{enumi}{3}
  \item \textbf{Osserva} cosa succede: Copilot creerà i file e potrebbe eseguire i test.
  \item \textbf{Prova il risultato:} \texttt{python calcolatrice.py}
\end{enumerate}
\end{praticabox}

\begin{checkpointbox}
Se la calcolatrice funziona e i test passano, il tuo ambiente è configurato correttamente. Hai appena creato il tuo primo programma in 0-code.
\end{checkpointbox}

% ---------------------------------------------------------
\section{Come Comunicare con l'IA: Le Basi}

\subsection{Le 5 regole della richiesta efficace}

\textbf{1.\ Sii specifico, non generico}
\begin{lstlisting}
# NO: "Crea un server"
# SI: "Crea un server Express.js su porta 3000 con un
#      endpoint GET /health che restituisca
#      { status: 'ok', timestamp: Date.now() }"
\end{lstlisting}

\textbf{2.\ Specifica la struttura dei file} --- mostra l'albero delle directory.

\textbf{3.\ Definisci i vincoli} --- cosa usare e cosa \textit{non} usare.

\textbf{4.\ Specifica il formato delle risposte} --- JSON, status code, errori.

\textbf{5.\ Dai il contesto del progetto} --- stack, architettura, dipendenze.

\begin{notebox}
Queste regole sono la versione semplificata della disciplina del \textbf{Context Engineering} che approfondiremo nel Capitolo~3 e applicheremo sistematicamente a partire dal Capitolo~5.
\end{notebox}

% ---------------------------------------------------------
\section{La Struttura di un Progetto 0-Code}

Ogni progetto che costruirai in questo libro avrà una struttura riconoscibile:

\begin{lstlisting}
mio-progetto/
|-- _CONTEXT.md          <- Il "contratto" con l'IA
|-- PROGRESS.md          <- Memoria persistente (progetti lunghi)
|-- .copilot/
|   |-- instructions.md  <- Regole specifiche del progetto
|   +-- skills/          <- Competenze specializzate (opzionali)
|       |-- react.md     <- SKILL frontend
|       +-- api-design.md <- SKILL API
|-- .gitignore
|-- README.md
|-- src/                 <- Codice sorgente (generato dall'IA)
|-- tests/               <- Test (generati dall'IA)
+-- package.json / requirements.txt / pubspec.yaml
\end{lstlisting}

I file evidenziati --- \texttt{\_CONTEXT.md}, \texttt{.copilot/}, \texttt{PROGRESS.md} --- sono gli unici file che \textbf{tu} scrivi manualmente. Tutto il resto è generato dall'IA sotto la tua supervisione.

\begin{quote}
Questo è il centro di gravità del paradigma 0-code: \textbf{tu scrivi il contesto, l'IA scrive il codice}.
\end{quote}

% ---------------------------------------------------------
\section{Configurazioni Consigliate per VS Code}

\begin{praticabox}[title={Pratica --- Settings consigliati}]
Apri le impostazioni di VS~Code (Command Palette $\rightarrow$ ``Open Settings (JSON)'') e aggiungi:
\begin{lstlisting}[language=json]
{
  "editor.formatOnSave": true,
  "editor.defaultFormatter": "esbenp.prettier-vscode",
  "editor.minimap.enabled": false,
  "files.autoSave": "onFocusChange",
  "terminal.integrated.defaultProfile.windows": "PowerShell",
  "github.copilot.chat.agent.enabled": true
}
\end{lstlisting}
\end{praticabox}

% ---------------------------------------------------------
\section*{Riepilogo}

\begin{center}
\begin{tabular}{l c}
\toprule
\textbf{Cosa hai fatto} & \textbf{Stato} \\
\midrule
Installato VS~Code & \checkmark \\
Installato GitHub Copilot & \checkmark \\
Scelto il modello IA (Claude Sonnet 4 raccomandato) & \checkmark \\
Generato il primo programma Python & \checkmark \\
Compreso le 5 regole della richiesta efficace & \checkmark \\
Configurato l'ambiente di lavoro & \checkmark \\
\bottomrule
\end{tabular}
\end{center}
